\section{先定义问题：这里的“推理”到底指什么}

关于“大语言模型是否真正会推理”的争论，经常把哲学、认知科学、形式逻辑和产品能力混在一起。Denny Zhou 在开场主动缩小问题：本讲只研究一种可操作、可测量的对象——模型在输入与最终答案之间生成的\textbf{中间 token 序列}。这个定义不要求轨迹像人脑，也不保证轨迹忠实解释模型内部计算；它只问：增加可见的中间生成步骤，是否能扩大模型可完成的计算、改变训练信号，并改善最终答案的选择。

\lecturefigure{slide-01.jpg}{课程标题：LLM Reasoning}{1}

\teachervoice{00:00:52--00:02:30，老师明确说自己不参加“LLM 到底算不算推理”的无定义争论。本讲把 reasoning 固定为 input 与 output 之间的 intermediate tokens；后续所有理论、解码和训练方法都围绕这个操作性定义展开。}

\subsection{中间 token 是计算接口，不是人格证明}

令输入为 $x$，中间推理序列为 $r=(r_1,\ldots,r_T)$，最终答案为 $a=(a_1,\ldots,a_M)$。自回归模型给整个响应分配概率：
\[
p_\theta(r,a\mid x)
=\prod_{t=1}^{T}p_\theta(r_t\mid x,r_{<t})
 \prod_{j=1}^{M}p_\theta(a_j\mid x,r,a_{<j}).
\]
其中 $\theta$ 是模型参数，$T$ 是中间 token 数，$M$ 是答案 token 数，$r_{<t}$ 与 $a_{<j}$ 分别表示已经生成的前缀。关键变化不是“模型突然拥有内心独白”，而是原本一次前向式的答案生成，被展开成 $T+M$ 个条件计算步骤；每个新 token 都能读取此前结果，并把临时状态写回上下文。

\lecturefigure{slide-02.jpg}{本讲的操作性定义：输入与输出之间的中间步骤}{2}

\begin{importantbox}{操作性定义}
本讲中的 LLM reasoning 指模型在最终答案之前生成的中间 token。它是一种\textbf{可观测计算轨迹}：可用于继续计算、训练、筛选和聚合；但“可读”不等于“忠实”，“像人”不等于“机制相同”。
\end{importantbox}

\begin{warningbox}{最常见的概念偷换}
看到自然语言步骤后，容易把“轨迹能帮助得到正确答案”偷换成“轨迹完整揭示了模型内部因果过程”。本讲支持前者，不自动推出后者。评价轨迹至少要分开看最终正确性、步骤有效性、可验证性与解释忠实性。
\end{warningbox}

\subsection{最后一个字母：为什么要设计不会被捷径破解的任务}

“artificial intelligence”的最后字母拼接为 “le”。直接答案只有两个字符，看不出模型如何获得它；分步版本则先抽取 artificial 的末字母 $l$，再抽取 intelligence 的末字母 $e$，最后拼接。这个小任务的教学价值不是难度，而是把一个串行过程拆成明确的临时状态，让研究者能区分“记住了常见模式”与“按步骤执行了转换”。

\lecturefigure{slide-03.jpg}{最后字母拼接：直接答案与中间步骤的对比}{3}

\teachervoice{00:03:18--00:05:10，老师解释最初试过“首字母拼接”，但网页中大量缩写让旧模型也能轻松答对，因此改用末字母。实践经验：合成任务必须检查训练语料中的捷径，否则 benchmark 测到的可能是共现记忆而不是目标算法。}

2017 年的 rationale generation 工作把自然语言中间步骤用于代数应用题；更早的神经符号研究常用程序、逻辑式或搜索状态表达中间过程。这里的历史转折不在于“第一次有中间状态”，而在于把状态放进与输入输出相同的语言序列中，从而可以直接用序列模型生成和学习。

\lecturefigure{slide-04.jpg}{Rationale generation：以自然语言承载中间计算}{4}

\begin{knowledgebox}{术语消化：rationale、scratchpad 与 Chain-of-Thought}
\textbf{rationale} 强调解释或解题依据；\textbf{scratchpad} 强调用可写的中间空间完成计算；\textbf{Chain-of-Thought（CoT）} 通常指连续的语言化步骤。三者在工程上可能使用相似 token 序列，但研究目标不同：解释、计算工作区与推理能力不应只靠名称混为一谈。
\end{knowledgebox}

\subsection{为什么中间 token 能增加可计算性}

官方讲义给出一个理论方向：若某问题可由规模为 $T$ 的 Boolean circuit 求解，那么固定规模的 Transformer 可以通过生成 $O(T)$ 个中间 token 来逐步模拟该计算；若强迫模型直接输出最终答案，则可能需要很深的网络，甚至无法在给定深度下表达。直觉上，网络深度提供“参数内的串行层数”，生成长度提供“推理时的串行步数”。

\lecturefigure{slide-05.jpg}{中间 token 将串行计算从网络深度转移到生成长度}{5}

可把总计算预算粗略写成
\[
C_{\text{total}}\approx C_{\text{prefill}}+T\,C_{\text{decode-step}}.
\]
其中 $C_{\text{prefill}}$ 是处理输入的成本，$C_{\text{decode-step}}$ 是生成一个新 token 的边际成本，$T$ 是中间步骤长度。更长轨迹提供更多串行计算，但同时增加延迟、KV cache 占用和错误累积，因此“允许更多 token”不是免费提升，而是把能力问题转成预算与选择问题。

\begin{importantbox}{计算视角的核心结论}
中间 token 的价值首先是\textbf{增加推理时可用的串行计算深度}。它解释了为什么同一参数模型在允许展开步骤后可能解决直接回答时失败的问题，也解释了为什么 test-time compute 必须明确“到底扩展了长度、样本数、搜索宽度还是工具调用”。
\end{importantbox}

\subsection{本章小结}

本章把 reasoning 固定为中间 token，并用概率分解、合成任务与计算深度解释其作用。这个定义刻意避免人类心智类比：后续真正要解决的是如何找到好轨迹、如何训练其分布、如何聚合多个轨迹，以及如何证明最终答案可靠。

